% titel: Blatt 7, A5: Eisberg
% skript: Physik I, Kapitel 8.1.3, Seiten 133--134
% stand: Übungsblatt 7, Aufgabe A5

Ein Eisberg schwimmt: Ein Teil ragt über, der grösste Teil liegt unter Wasser.
Gefragt ist nach dem Zusammenhang zwischen eingetauchtem Volumen, Dichte und
so weiter. Das Werkzeug ist der Auftrieb in einer ruhenden Flüssigkeit.

\begin{werkzeugkasten}[Dichte und Auftrieb]
Die Dichte verbindet Masse und Volumen eines homogenen Stoffes über
$m=\rho V$. Ein tauchfähiger Körper erfährt eine nach oben gerichtete
Auftriebskraft, die dem Gewicht des von ihm verdrängten Fluids entspricht:
\begin{equation}
  F_\mathrm A=\rho_\mathrm F\, g\, V_{\text{verdrängt}}.
\end{equation}
Entscheidend ist das \emph{verdrängte} Volumen, das beim Schwimmen nur der
eingetauchte Teil ist. Im Gleichgewicht heben sich Auftrieb und Gewicht des
Körpers ($mg=\rho_\mathrm K g V$) auf.

\textit{Skript:} Kap. 8.1.3, S. 133--134, Gl. (8.5).
\end{werkzeugkasten}

\begin{werkzeugkasten}[Dichte des Körpers und eingetauchter Anteil]
Im Gleichgewicht gilt $\rho_\mathrm K V=\rho_\mathrm F V_{\text{verdrängt}}$.
Je dichter der Körper im Vergleich zum Fluid, desto grösser der eingetauchte
Anteil; ein Körper mit $\rho_\mathrm K<\rho_\mathrm F$ schwimmt, Teile ragen
heraus. So lässt sich aus den Dichten sofort der Schwimm- und der
eingetauchte Anteil abschätzen, ohne die Form zu kennen.

\textit{Skript:} Kap. 8.1.3, S. 133--134.
\end{werkzeugkasten}

\begin{achtung}[Darauf achten]
\begin{itemize}
  \item Im Gleichgewicht sind es zwei verschiedene Kräfte (Auftrieb und
        Gewicht), die sich ausgleichen; die Dichten von Körper und Fluid nicht
        verwechseln.
  \item Das verdrängte Volumen ist beim Schwimmen nicht das ganze
        Körpervolumen.
\end{itemize}
\end{achtung}
